% !TEX root = ../main.tex
% SECTION 0: WICHTIGE FORMELN
% ============================================================

\StartChapterOnNewColumn[chap:wichtige]{Wichtige Formeln}

\begin{runintext}
Kompakter Querschnitt zentraler Beziehungen aus den Kapiteln 1--11.
\end{runintext}

\begin{runintext}
Grundlagen der Elektrostatik: Kraft zwischen zwei Ladungen und das von einer Punktladung erzeugte elektrische Feld.
\end{runintext}
\begin{formulabox}
\[ \vec{F} = \frac{1}{4\pi\varepsilon_0}\,\frac{q_1 q_2}{r^2}\,\hat{r} \quad \text{Coulomb} \]
\formulasep
\[ \vec{E} = \frac{1}{4\pi\varepsilon_0}\,\frac{q}{r^2}\,\hat{r} \quad \text{E-Feld Punktladung} \]
\formulanote{Zweck: Kraft zwischen Ladungen; E-Feld einer Punktladung.}
\end{formulabox}

\begin{runintext}
Zusammenhang zwischen Ladung, elektrischem Fluss und Potenzial. Wichtig für Aufgaben mit Symmetrie oder Energieberechnung.
\end{runintext}
\begin{formulabox}
\[ \oint \vec{E}\cdot\mathrm{d}\vec{A} = \frac{q_{\mathrm{in}}}{\varepsilon_0} \quad \text{Gauss E} \]
\formulasep
\[ \varphi = \frac{1}{4\pi\varepsilon_0}\,\frac{q}{r}, \quad \vec{E} = -\nabla\varphi \quad \text{Potenzial} \]
\formulanote{Zweck: Gauss (Fluss $\propto$ Ladung); Potenzial und E $= -\nabla\varphi$.}
\end{formulabox}

\begin{runintext}
\textbf{Dünne Kugelschale} (Radius $R$, Ladung $Q$): Im Inneren ist das E-Feld null (Faraday-Käfig) und das Potenzial ortsunabhängig konstant. Aussen wirkt die Schale wie eine Punktladung im Zentrum. Wichtig für Hohlleiter, geerdete Kugeln und Influenz-Aufgaben.
\end{runintext}
\begin{formulabox}
\[
\text{Schale:} \begin{cases}
E = 0, \hfill & r < R \\
\varphi = \frac{1}{4\pi\varepsilon_0}\frac{Q}{R}, \hfill & r < R \\[1ex]
E = \frac{1}{4\pi\varepsilon_0}\frac{Q}{r^2}, \hfill & r > R \\
\varphi = \frac{1}{4\pi\varepsilon_0}\frac{Q}{r}, \hfill & r > R
\end{cases}
\]
\formulanote{Zweck: E-Feld \& Potenzial einer Kugelschale. \ZSFkeyword{Erden}: Zwingt Potenzial auf $\varphi = 0$ ($\infty$-Speicher). Bei \ZSFkeyword{Influenz} (geerdet): zugewandte Seite lädt sich entgegengesetzt, abgewandte bleibt neutral (Ladungsabfluss).}
\end{formulabox}

\begin{runintext}
Superposition: Gesamtpotenzial einer gleichmässig geladenen \textbf{Vollkugel} ($R_1$, $Q_1$) innerhalb zweier konzentrischer \textbf{Kugelschalen} ($R_2, R_3$). Oft als Prüfungstrick geprüft. Mit $k = \frac{1}{4\pi\varepsilon_0}$:
\end{runintext}
\begin{formulabox}
\[
\varphi(r) = k \begin{cases}
 \begin{aligned}&\tfrac{Q_1}{2R_1^3}(3R_1^2 - r^2)\\ &+ \tfrac{Q_2}{R_2} + \tfrac{Q_3}{R_3}\end{aligned}, & r < R_1 \\[2ex]
 \frac{Q_1}{r} + \frac{Q_2}{R_2} + \frac{Q_3}{R_3}, & R_1 \le r < R_2 \\[1ex]
 \frac{Q_1 + Q_2}{r} + \frac{Q_3}{R_3}, & R_2 \le r < R_3 \\[1ex]
 \frac{Q_1 + Q_2 + Q_3}{r}, & r \ge R_3
\end{cases}
\]
\formulanote{Zweck: Gesamtpotenzial addiert sich aus den inneren und äusseren Schalenbeiträgen (innen = konstant, aussen = $1/r$).}
\end{formulabox}

\begin{runintext}
Kapazität und Energiespeicherung, insbesondere für den idealisierten Plattenkondensator.
\end{runintext}
\begin{formulabox}
\[ C = \frac{q}{U} \quad \text{Kapazität} \]
\formulasep
\[ C = \frac{\varepsilon_0 A}{d} \quad \text{Platte}, \quad E_{\mathrm{el}} = \frac{1}{2}CU^2 \]
\formulanote{Zweck: Kapazität; Plattenkondensator; gespeicherte Energie.}
\end{formulabox}

\begin{runintext}
Grundgleichungen für Stromkreise: Ohmsches Gesetz, Kirchhoff'sche Regeln für Maschen/Knoten und Zeitkonstante für das Entladen.
\end{runintext}
\begin{formulabox}
\[
\begin{gathered}
U = R\,I, \quad R = \rho\,\ell/A \quad \text{Ohm} \\
\sum U = 0 \;\text{(Maschen)},\; \sum I_{\mathrm{in}} = \sum I_{\mathrm{out}} \;\text{(Knoten)} \\
\tau = RC
\end{gathered}
\]
\formulanote{Zweck: Ohm; Kirchhoff; RC-Zeitkonstante.}
\end{formulabox}

\begin{runintext}
Magnetostatik: Kraft auf bewegte Ladungen, Magnetfeld von Strömen und das Gesetz von Ampère.
\end{runintext}
\begin{formulabox}
\[ \vec{F} = q\,\vec{v}\times\vec{B} \quad \text{Lorentz} \]
\formulasep
\[ \mathrm{d}\vec{B} = \frac{\mu_0}{4\pi}\,\frac{I\,\mathrm{d}\vec{\ell}\times\hat{r}}{r^2} \quad \text{Biot-Savart} \]
\formulasep
\[ \oint \vec{B}\cdot\mathrm{d}\vec{\ell} = \mu_0 I \quad \text{Ampère} \]
\formulanote{Zweck: Lorentzkraft; B-Feld Stromelement; Ampère.}
\end{formulabox}

\begin{runintext}
\ZSFkeyword{Winkelparametrisierung}: Geniales Rezept für Integrale über Linienleiter ($\vec{E}, \vec{B}$). Wandelt die unschöne Integration über den Ort $x$ in eine simple trigonometrische Integration über den Winkel $\vartheta$ um.
\end{runintext}
\begin{procedure}[Winkelparametrisierung]
  \ProcStep{Substitution}{%
    $h$ = Lotabstand, $\vartheta$ = Winkel relativ zum Lot.
    \[ \begin{aligned}
        x &= h \cdot \tan\vartheta \quad\implies\quad \mathrm{d}x = \frac{h}{\cos^2\vartheta}\,\mathrm{d}\vartheta \\[1ex]
        r &= \frac{h}{\cos\vartheta} \quad\implies\quad \frac{\mathrm{d}x}{r^2} = \frac{1}{h}\,\mathrm{d}\vartheta
    \end{aligned} \]
  }
  \ProcStep{E-Feld einer Linienladung $\lambda$}{%
    \[ \begin{aligned}
        E_{\parallel} &= \frac{\lambda}{4\pi\varepsilon_0 h} \int_{\vartheta_1}^{\vartheta_2} \sin\vartheta\,\mathrm{d}\vartheta = \frac{\lambda}{4\pi\varepsilon_0 h} \Bigl[-\cos\vartheta\Bigr]_{\vartheta_1}^{\vartheta_2} \\[2ex]
        E_{\perp} &= \frac{\lambda}{4\pi\varepsilon_0 h} \int_{\vartheta_1}^{\vartheta_2} \cos\vartheta\,\mathrm{d}\vartheta = \frac{\lambda}{4\pi\varepsilon_0 h} \Bigl[\sin\vartheta\Bigr]_{\vartheta_1}^{\vartheta_2}
    \end{aligned} \]
  }
  \ProcStep{B-Feld eines geraden Stromleiters $I$}{%
    \[
        B = \frac{\mu_0 I}{4\pi h} \int_{\vartheta_1}^{\vartheta_2} \cos\vartheta\,\mathrm{d}\vartheta = \frac{\mu_0 I}{4\pi h} \Bigl[\sin\vartheta\Bigr]_{\vartheta_1}^{\vartheta_2}
    \]
    \formulanote{Unendlicher Leiter ($\vartheta \in [-\pi/2, \pi/2]$): $E_{\perp} = \frac{\lambda}{2\pi\varepsilon_0 h}$, $B = \frac{\mu_0 I}{2\pi h}$. Komponente $E_{\parallel}$ hebt sich auf.}
  }
\end{procedure}

\begin{runintext}
Elektromagnetische Induktion nach Faraday und Energie im Magnetfeld einer Spule.
\end{runintext}
\begin{formulabox}
\[ U_{\mathrm{ind}} = -\frac{\mathrm{d}\Phi_{\mathrm{mag}}}{\mathrm{d}t} \quad \text{Faraday} \]
\formulasep
\[ n\,\Phi_{\mathrm{mag}} = L\,I, \quad E_{\mathrm{mag}} = \frac{1}{2}LI^2 \]
\formulanote{Zweck: Faraday; Selbstinduktion; magnetische Energie.}
\end{formulabox}

\begin{runintext}
Wechselstrom: Widerstände von Spule/Kondensator, Eigenfrequenz eines Schwingkreises und Spannungsübersetzung.
\end{runintext}
\begin{formulabox}
\[ \omega = 2\pi\nu = 2\pi f = \frac{2\pi}{T} \quad \text{Frequenz/Kreisfrequenz} \]
\formulasep
\[ X_L = \omega L, \quad X_C = 1/(\omega C), \quad \omega_0 = 1/\sqrt{LC} \quad \text{LC} \]
\formulasep
\[ U_2/U_1 = n_2/n_1 \quad \text{Transformator} \]
\formulanote{Zweck: Umrechnung Frequenz ($\nu, f$) / Kreisfrequenz $\omega$ (Achtung: Frequenz $\nu \neq v$); Blindwiderstände; LC-Eigenfrequenz; Transformator-Übersetzung.}
\end{formulabox}

\begin{runintext}
Elektromagnetische Wellen: Ausbreitungsgeschwindigkeit, Verhältnis von E zu B und Energietransport.
\end{runintext}
\begin{formulabox}
\[ c = 1/\sqrt{\mu_0\varepsilon_0}, \quad E = c\,B \quad \text{EM-Welle} \]
\formulasep
\[ \vec{S} = \vec{E}\times\vec{B}/\mu_0 \quad \text{Poynting} \]
\formulanote{Zweck: Lichtgeschwindigkeit; E/B in Welle; Poynting-Vektor.}
\end{formulabox}

\begin{runintext}
Optik: Brechungsgesetz, Abbildungsgleichung für Linsen/Spiegel und Beugung am Gitter.
\end{runintext}
\begin{formulabox}
\[ n_1\sin\theta_1 = n_2\sin\theta_2 \quad \text{Snellius} \]
\formulasep
\[ \frac{1}{g}+\frac{1}{b} = \frac{1}{f} \quad \text{Spiegel/Linse} \]
\formulasep
\[ g\,\sin\theta_m = m\,\lambda \quad \text{Gitter} \]
\formulanote{Zweck: Brechung; Abbildung Spiegel/Linse; Gitter-Maxima.}
\end{formulabox}

% ============================================================
